\documentclass[12pt]{article}

\usepackage{bbm}
\usepackage{fancyhdr}
\usepackage{titling}
\usepackage{enumitem}
\usepackage{amsmath}
\usepackage{dsfont}
\usepackage{amsfonts}
\usepackage{amssymb}
\usepackage{amsthm}
\usepackage{geometry}
\usepackage{tikz}
\usepackage{algorithm}
\usepackage{algpseudocode}
\geometry{verbose,tmargin=1in,bmargin=1in,lmargin=1in,rmargin=1in}
\renewcommand{\vec}[1]{\mathbf{#1}}
\newcommand{\dec}{\text{dec}}
\newcommand{\divi}{\mathbin{\textnormal{div}}}
\title{Homework 1}
\author{Satvik Prasad}
\date{}
\newtheorem{lemma}{Lemma}[section]
\fancypagestyle{plain}{%
  \fancyhf{}
  \fancyhead[L]{\theauthor}
  \fancyhead[R]{\thedate}
}
\pagestyle{plain}
\makeatletter
\def\@maketitle{%
  \newpage
  \null
  \vskip 1em%
  \begin{center}%
    \let \footnote \thanks
    {\LARGE \@title \par}%
    \vskip 1em%
  \end{center}%
  \par
\vskip 1em}
\makeatother
\usepackage{cmbright}

\newcommand{\mat}[1]{\mathbf{#1}}

% ---------- Probability theory ----------
% Core operators (optional argument-free; use as \E{X}, \Var{X}, \Cov{X}{Y})
\DeclareMathOperator{\EOp}{\mathbb{E}}
\DeclareMathOperator{\VarOp}{Var}
\DeclareMathOperator{\CovOp}{Cov}
\DeclareMathOperator{\CorrOp}{Corr}

\renewcommand{\P}[1]{\mathbb{P}\left(#1\right)}        % \P{A}
\newcommand{\E}[1]{\EOp\left[#1\right]}                % \E{X}
\newcommand{\Var}[1]{\VarOp\left(#1\right)}            % \Var{X}
\newcommand{\Cov}[2]{\CovOp\left(#1,\,#2\right)}       % \Cov{X}{Y}
\newcommand{\Corr}[2]{\CorrOp\left(#1,\,#2\right)}     % \Corr{X}{Y}
\newcommand{\SD}[1]{\operatorname{SD}\left(#1\right)}  % \SD{X}

% Subscripted versions, e.g. \Esub{\theta}{X}, \Psub{\theta}{A}
\newcommand{\Esub}[2]{\EOp_{#1}\left[#2\right]}
\newcommand{\Psub}[2]{\mathbb{P}_{#1}\left(#2\right)}

% Conditioning, indicators, independence
\newcommand{\given}{\,\middle|\,}                      % \E{X \given Y}
\newcommand{\ind}[1]{\mathbbm{1}\left\{#1\right\}}     % \ind{X > 0}
\newcommand{\indep}{\perp\!\!\!\perp}                  % X \indep Y
\newcommand{\iid}{\overset{\text{iid}}{\sim}}          % X_i \iid \Normal(0,1)

% Convergence
\newcommand{\convp}{\xrightarrow{\;p\;}}               % in probability
\newcommand{\convd}{\xrightarrow{\;d\;}}               % in distribution
\newcommand{\convas}{\xrightarrow{\text{a.s.}}}        % almost surely
\newcommand{\convLp}[1]{\xrightarrow{L^{#1}}}          % in L^p

% Probability spaces
\newcommand{\F}{\mathcal{F}}                           % sigma-algebra
\newcommand{\B}{\mathcal{B}}                           % Borel sigma-algebra
\newcommand{\pspace}{(\Omega, \F, \mathbb{P})}

% Common distributions
\newcommand{\Normal}{\mathcal{N}}
\newcommand{\Bern}{\operatorname{Bernoulli}}
\newcommand{\Bin}{\operatorname{Binomial}}
\newcommand{\Pois}{\operatorname{Poisson}}
\newcommand{\Unif}{\operatorname{Uniform}}
\newcommand{\Exp}{\operatorname{Exp}}
\newcommand{\Geom}{\operatorname{Geometric}}
\newcommand{\Gam}{\operatorname{Gamma}}
\newcommand{\Beta}{\operatorname{Beta}}

% Number sets
\newcommand{\R}{\mathbb{R}}
\newcommand{\N}{\mathbb{N}}
\newcommand{\Z}{\mathbb{Z}}


\begin{document}
\subsection{Green Circle}
\begin{enumerate}
  \item [1] We want to count the number of ordered pairs $(A, B)$ of
    disjoint subsets of $[n].$ Since $A, B.$ Consider any such
    ordered pair $(A, B).$ Define the mapping $f: D \to (0, 1,
    2)^{n},$ such that
    $$
    f(A, B)_{i} =
    \begin{cases}
      0 & i \in A  \\
      1 & i \in B \\
      2 & \text{otherwise}
    \end{cases}.
    $$
    This mapping is a bijection, and hence we can count $|(0, 1, 2)^n| = 3^n.$

  \item [4] We want to determine the number of ways to partition $2n$
    distinct elements into $n$ unordered pairs.

    First, consider counting such a quantity where each pair is
    distinct. I.e., define buckets $1, 2, \ldots, n$ that can contain
    maximum $2$ elements. There are
    $$
    \dfrac{(2n)!}{(2!)^{n}} = \dfrac{(2n)!}{2^n}
    $$
    ways to perform this permutation, since each permutation of $1,
    \ldots, 2n$ corresponds with $2^n$ selections of unordered pairs
    where pairs are distinct.

    Then, since pairs are indistinct, there are $n!$ cases with
    distinct pairs that correspond to a single case with indistinct
    pairs (partition). Thus, the total number of ways is
    $$
    \dfrac{(2n)!}{n! 2^n}.
    $$

  \item [5] We want a combinatorial proof for
    $$
    K :=\sum_{k =1 }^{n} k {n \choose k}^2 = n {2n - 1 \choose n - 1}.
    $$
    when $n \geq 1.$

    Consider $n$ elements from group A, and $n - 1$ elements from
    group B, with a leader from group B. Naively, we can do this by
    selecting the leader in $n$ ways, selecting the remaining $n -1$
    elements in ${2n -1 \choose n - 1}$ ways, yielding
    $$
    K = n {2 n -1 \choose n- 1}.
    $$

    Alternatively, we can choose an arbitrary $1 \leq k \leq n$
    elements from $B,$ and then choose $n - k$ elements to not be
    chosen from $A,$ and then choose which element is the leader from $B$ in $$
    k {n \choose k}^2
    $$
    ways. Summing over all possible $k$, we therefore get
    $$
    \sum_{k =1 }^{n} k {n \choose k}^2 = n {2n -1 \choose n - 1}.
    $$
\end{enumerate}

\subsection{Blue Square}
\begin{enumerate}
  \item [7] We want to find the probability that a permutation of
    $[n]$ has exactly 2 cycles.

    Consider writing an arbitrary cyclic permutation of $[n]$ with
    two cycles in cyclic notation, where the length of the first
    cycle is $\leq$ the length of the second cycle.

    For the following, assume $n$ is odd. Let $r = (n - 1)/2.$

    Then, if $n$ is odd every cycle can be written in
    $$
    \sum_{k = 1}^{r}{n \choose k} (k-1)! (n- k -1)! = \sum_{k =1
    }^{r} \dfrac{n!}{k (n - k)}
    $$
    ways.

    There are $n!$ total permutations, meaning we have a probability
    $$
    \sum_{k = 1}^{r} \dfrac{1}{k(n-k)}
    $$
    for a permutation to have exactly $2$ cycles.

    Since $$
    \dfrac{1}{k(n-k)} = \dfrac{1}{n}\left(\dfrac{1}{k} + \dfrac{1}{n -k}\right),
    $$
    we can write the above as
    \begin{align*}
      \dfrac{1}{n}\sum_{k =1 }^{r}  \dfrac{1}{k} + \dfrac{1}{n - k} +
      \dfrac{1}{n} \left( \sum_{k = 1}^{r} \dfrac{1}{k}
        + \sum_{k = \underbrace{n -r }_{r + 1}}^{n -1} \dfrac{1}{ k}
      \right)
      =  \dfrac{1}{n}\left(1 + \dfrac{1}{2} + \ldots + \dfrac{1}{n-1}\right)
    \end{align*}

    In the case where $n$ is even, the $k = n/2$ case should not have
    a double count, meaning our sum is
    $$
    \underbrace{\dfrac{1}{2}{n \choose n/2}(n/2 - 1)!^2}_{=
    \frac{1}{2}(\frac{n!}{n^2})}   + \sum_{k = 1}^{n /2 - 1}
    \dfrac{n!}{k(n-k)} = \dfrac{1}{n} \cdot \dfrac{n!}{2n} + \sum_{k
    =1 }^{n/2 - 1} \dfrac{n!}{k(n-k)}.
    $$
    Hence, our probability (dividing by $n!$) is
    $$
    \dfrac{1}{n}\left(\dfrac{1}{2n}\right) +\dfrac{1}n{} \sum_{k =
    1}^{n/2 - 1} \dfrac{1}{k} + \dfrac{1}{n}\sum_{k = n /2  +
    1}^{n-1} \dfrac{1}{k},
    $$
    which is also
    $$
    \dfrac{1}{n}\left(
      1 + \dfrac{1}{2} + \ldots + \dfrac{1}{n}
    \right)
    $$
    as required.

  \item [9] We prove that every permutation $\sigma$ of $[n]$ can be
    written $\sigma = \sigma_1 \sigma_2$ where $\sigma_1, \sigma_2$
    are involutions, and compositions are performed right to left.

    Consider the cycle notation of an arbitrary permutation $\tau.$
    For $\tau^2 = \mathds{1},$ the element following $a$ in some
    cycle must be also followed by $a$ (up to some permutation), or
    $(a)$ msut be in it's own cycle. In other words, a permutation
    satisfies $\tau^2 = \mathds 1$ if and only if it has cycles of at
    most length $2,$ i.e composed of only disjoint swaps.

    First, we want to show that we can always write some arbitrary
    cycle as a product of two involutions. Consider an arbitrary cycle
    $$
    \sigma = (a_1 \ a_2 \ \ldots \ a_k),
    $$
    for some integer $k.$ Assume we have $k \equiv 0 \pmod 2.$ Then, a cycle
    $$
    (a_1 \ a_2 \ \ldots \ a_k)
    $$
    can be written as $\alpha \beta,$ where
    $$
    \beta = (a_1 \ a_k) (a_2 \ a_{k -1 }) \ldots (a_{k/2} a_{k/2 + 1})
    $$
    and
    $$
    \alpha = (a_k \ a_2) (a_{k - 1} a_3) \ldots (a_{k/2 + 2} \ a_{k/2}).
    $$
    Similarly, if we have $k \equiv 1 \pmod 2,$ and $k' = (k -1)/2$
    then a cycle can be decomposed as
    $$
    \beta = (a_1 \ a_k) (a_2 \ a_{k-1}) \ldots (a_r \ a_{k + 1 - r}),,,,
    $$
    and
    $$
    \alpha = (a_k \ a_2) (a_{k-1} \ a_3) \ldots (a_{k + 1- r} \ a_{r + 1}).
    $$

    In either case, both $\alpha^2 = \mathds 1$ and $\beta^2 =
    \mathds 1,$ meaning we can write an arbitrary cycle $c = c_1 c_2$
    where $c_1, c_2$ involutions.

    Then, decompose $\sigma$ into a product of arbitrary (disjoint)
    cycles $\sigma_j$ (if cycles aren't disjoint, combine them until
    they are). By the above, we can write
    $$
    \sigma = \prod_{j} \alpha_j \beta_j.
    $$
    Since $\alpha_j, \beta_i$ for $j \neq i$ are disjoint, they
    commute. Thus, by iteratively commuting $\beta_i$ all the way to
    the right for decreasing $i,$ we can write
    $$
    \sigma = \underbrace{\prod_{j} \alpha_j}_{\sigma_1} \times
    \underbrace{\prod_{j} \beta_j.}_{\sigma_2}
    $$
    Since $\beta_i, \beta_j$ are disjoint for $i \neq j,$ and so are
    $\alpha_i, \alpha_j,$ $\sigma$ can be written as $\sigma =
    \sigma_1 \sigma_2$ for involutions $\sigma_1, \sigma_2.$
\end{enumerate}

\subsection{Black Diamond}
\begin{enumerate}
  \item [10] We want to count the number of minimal selfish subsets of $[n].$

    Consider any minimal selfish subset $A \subseteq [n]$ with $|A| =
    k.$ Then, $k \in A$ since $A$ is selfish. Now, for the sake of
    contradiction assume $r \in A$ for $1 \leq r < k.$ Then, we can
    construct a subset $A' \subseteq A$ with $r \in A,$ and $r - 1 <
    k - 1$ other arbitrary elements from $A$ in $A',$ meaning $A'$ is
    also selfish, and hence $A$ cannot be minimal selfish. Thus, for
    minimal selfish $A \subseteq [n],$ we must have $a \geq |A|$ for
    all $a \in A.$

    Alternatively, consider any subset $A \subseteq [n]$ with $|A| =
    k,$ $k \in A,$ and $r \geq k$ for all $r \in A.$ Then, any proper
    subset $A' \subset A$ has cardinality $|A'| < k,$ meaning no
    element $r \in A'$ can have $r = |A'|,$ since $r \geq k > |A'|,$
    and thus $A$ is minimal selfish.

    Hence, a set $A \subseteq [n]$ is minimal and selfish if and only
    if for all $r \in A,$ $r \geq |A|,$ and $|A| \in A.$ Then, we can
    count this quantity by enumerating over $|A|$:
    \begin{align*}
      \sum_{k = 1}^{n} \underbrace{{n - k \choose k - 1}}_\text{ways
      to choose $k - 1$ other elements $> k$},
    \end{align*}
    as required.
\end{enumerate}
\subsection{Reflection}
Initially, for Blue Square (7), didn't consider the double-counting
for even cases correctly, and got issues with my summation tricks.
Repaired it by realizing every unordered case is counted twice
(appearing in either the left cycle or the right cycle), and then
appropriately dividing by 2.
\end{document}
